% MỞ ĐẦU
% Quy định trình bày (mục 2): nêu tính cần thiết, ý nghĩa khoa học và thực tiễn,
% đối tượng và phương pháp nghiên cứu, nội dung nghiên cứu.
% Rubric: Chủ đề, Đóng góp chính, Bố cục.
\chuongkhongso{Mở đầu}

Dự báo sự kiện theo chuỗi thời gian (time-series event forecasting) --- chẳng hạn dự báo một
cơn động kinh sắp xảy ra, hay dự báo một biến đổi mạnh trong một hệ động lực --- luôn đối mặt
với một đánh đổi cố hữu giữa tính hữu dụng và độ chính xác. Dự báo càng sớm thì càng hữu dụng
về mặt ứng dụng (có thời gian để phản ứng) nhưng càng khó và kém chính xác, vì tín hiệu báo
trước còn mơ hồ; dự báo càng muộn, gần sát thời điểm sự kiện xảy ra, thì càng dễ và chính xác
nhưng lại mất giá trị thực tiễn, vì không còn đủ thời gian để hành động. Phần lớn các mô hình
dự báo truyền thống huấn luyện một mô hình duy nhất để tự giải quyết đánh đổi này, khiến mô
hình phải cân bằng ngầm giữa hai mục tiêu trái chiều mà không có cơ chế tách biệt rõ ràng.

Future-Guided Learning (FGL), phương pháp được đề xuất trong bài báo
``A predictive approach to enhance time-series forecasting'' của Gunasekaran và cộng sự,
công bố trên \textit{Nature Communications} năm 2025, giải quyết đánh đổi trên bằng một ý
tưởng đơn giản nhưng hiệu quả: tách bài toán thành hai mô hình riêng biệt --- một
\textit{teacher} được huấn luyện trên nhiệm vụ dễ hơn, dùng dữ liệu ở tương lai gần hơn so
với mô hình thứ hai; và một \textit{student} học nhiệm vụ khó thực sự cần giải quyết, dùng dữ
liệu ở quá khứ tương đối. Trong lúc huấn luyện, teacher dẫn dắt student qua cơ chế chưng cất
tri thức (knowledge distillation), giúp student học nhanh và chính xác hơn so với khi tự học
một mình.

Ý nghĩa khoa học của FGL nằm ở việc đây là một khung phương pháp tổng quát, áp dụng được cho
cả bài toán hồi quy (dự báo giá trị liên tục của chuỗi thời gian) và bài toán phân loại sự
kiện (dự báo một biến cố nhị phân), chỉ bằng cách thay đổi cách định nghĩa ``tương lai'' mà
teacher được nhìn thấy. Về mặt thực tiễn, các bài toán mà FGL hướng tới --- dự báo cơn động
kinh từ điện não đồ, dự báo chuỗi thời gian có cấu trúc thay đổi theo thời gian --- đều là các
bài toán có giá trị ứng dụng cao trong y tế và trong các lĩnh vực cần cảnh báo sớm.

Đối tượng nghiên cứu của báo cáo này là phương pháp FGL và mã nguồn công khai do tác giả công
bố (kho \texttt{SkyeGunasekaran/FutureGuidedLearning} trên GitHub). Phương pháp nghiên cứu là
đọc trực tiếp mã nguồn để hiểu cách phương pháp được hiện thực hóa, tái hiện lại thực nghiệm
trên hai trong ba bộ dữ liệu của bài báo --- chuỗi hỗn loạn Mackey--Glass (dự báo hồi quy) và
điện não đồ CHB-MIT (dự báo cơn động kinh) --- rồi đối chiếu kết quả tái hiện với số liệu công
bố. Báo cáo này tập trung vào việc tái hiện trung thực thực nghiệm của bài báo gốc; các cải
tiến tự đề xuất (so sánh kiến trúc hồi quy, LSTM hai chiều) được dành cho báo cáo cuối kỳ.

Đóng góp chính của báo cáo gồm hai phần. Thứ nhất, báo cáo tái hiện thành công thực nghiệm
Mackey--Glass, phát hiện và sửa tám lỗi cụ thể trong mã nguồn công khai của tác giả (từ các
tệp còn nguyên dấu xung đột merge khiến mã không chạy được, đến một khối tham số dòng lệnh bị
vô hiệu do nằm trong chuỗi chú thích), sau đó chạy lại trên nhiều seed để kiểm tra độ ổn định
của kết quả. Thứ hai, báo cáo phân tích chi tiết cơ chế FGL trên bài toán dự báo cơn động kinh
CHB-MIT và tái hiện đầy đủ cả ba giai đoạn trên ba bệnh nhân (1, 3, 5; sau khi chuyển từ máy cá
nhân sang Kaggle do hết bộ nhớ), cho thấy mức tăng AUC-ROC trung bình $+31,2\%$, thấp hơn mức
44,8\% bài báo công bố --- một kết quả không ủng hộ hoàn toàn tuyên bố của bài báo, được trình
bày trung thực cùng các phân tích khả năng về nguyên nhân,
thay vì bị bỏ qua hoặc chỉ chọn số liệu có lợi.

Phần còn lại của báo cáo được tổ chức như sau. Chương 1 phát biểu bài toán dự báo sự kiện mà
FGL hướng tới giải quyết. Chương 2 trình bày kiến thức nền tảng cần thiết để hiểu phương pháp
của bài báo gốc, gồm mạng hồi quy RNN, chưng cất tri thức và các khái niệm liên quan tới hai bộ
dữ liệu thực nghiệm. Chương 3 trình bày chi tiết phương pháp FGL và cách phương pháp này được
hiện thực hóa khác nhau trên từng bài toán. Chương 4 mô tả dữ liệu, công cụ và kịch bản thực
nghiệm. Chương 5 trình bày kết quả tái hiện và các vấn đề phát hiện được trong mã nguồn tác
giả. Chương 6 đánh giá so sánh kết quả tái hiện với kết quả công bố và với các hướng tiếp cận
liên quan. Cuối cùng là kết luận, tóm tắt kết quả đạt được và hướng phát triển.
